\section{Benchmark Protocol}

\label{sec:protocol}



The benchmark assesses the visual association aspect by determining, for a new crack observation, whether it matches one of the recorded defects or if it does not appear in the reference set. Since annotated identities act as stand-ins for permanent defect records, the measurement is of association accuracy rather than duplicate-ticket reduction or overall workflow performance.



\subsection{Task formulation and candidate records}



A wall means a single physical wall; an identity is the association between observations relating to the same crack (as explained in Sec.~\ref{sec:annotation}). A query involves a single crack instance; its gallery comprises other instances from photographs of the same wall, provided that the wall is known. Cross-wall search is not included in this approach. A query is regarded as known if its gallery contains the same identity, otherwise it is unknown. Each method assigns scores to the query–gallery pairs in order to carry out ranking, thresholding decisions, and identification with rejection. The unit used in pairs is a query–gallery pair; an image pair is the unit of registration or assignment. The Hungarian one-to-one constraint is applied only in the case of each image pair; more than one observation may refer to the same defect record.



The numbers of the tests and how common they are are presented in Table~\ref{tab:units} for the frozen benchmark snapshot (this covers 6 validation walls, namely walls numbered 01 to 06; 11 test walls, namely walls numbered 07 to 17; the walls labelled wall18 and wall19, which were added later, are omitted from all the figures). Registration is performed for all 301 pairs of test images that are identical in wall; 255 of these result in scored pairs with validity-mask cells (for further details, refer to Sec.~\ref{sec:coverage}).



\begin{table}[t]

\centering

\small

\caption{Counting units for every metric (frozen snapshot, $n_Q{=}305$ test).}

\label{tab:units}

\begin{tabular}{lrr}

\toprule

 & Validation & Test \\

\midrule

Walls / photographs / identities & 6 / 66 / 66 & 11 / 74 / 146 \\

Identities in $\geq 2$ photographs & 18 & 19 \\

Answerable queries & 439 & 305 \\

Positive query--gallery pairs & 26\,012 & 5\,722 \\

Valid query--gallery pairs & 89\,118 & 31\,822 \\

Same-wall image pairs & 394 & 301 \\

Scored image pairs (with valid cells) & 341 & 255 \\

\bottomrule

\end{tabular}

\end{table}



\subsection{Unchanged damage under different acquisition conditions}



The main evaluation involves carrying out real walk-around inspections on a per-wall basis, with the morphology regarded as essentially unchanged and any differences arising from variations in viewpoint, distance, scale, lighting, and visibility, together with the inclusion of same-wall distractors and unknown queries. The walls are not shared between the validation and test phases. Each method is used exactly as it is without any fine-tuning using \textsc{CrackID} on shared masks (Sec.~\ref{sec:method}); the thresholds are set using the validation data and then applied to the test data. In this way, within-visit correspondence is established, but neither cross-visit nor progression robustness is addressed.



\subsection{Controlled appearance change and damage progression}



A synthetic reanalysis separates brightness variations from those caused by structural growth by using ten photographs from five walls, along with one initial seed, and applying the same slight affine warp in each case. In condition B, the appearance is degraded (with blur, noise, and changes in contrast or brightness) even though the structure remains unchanged; in condition A, one end of the largest component is extended. The code specifies all parameters. The study compares OSNet with context to the \hybrid{} method (which combines geometric alignment with a fallback/blend based on the skeleton, $\tau{=}8.0$\, px), which differs from the main \ours{} reference. Because we use only one seed and five walls, the results are preliminary checks, not longitudinal evidence.



\subsection{Decision rules and handling of unavailable scores}



When a pair meets the threshold, it is accepted; at the query level, if the threshold is met then only the top entry from the gallery is accepted, otherwise no acceptance occurs; in the case of assignment, the Hungarian algorithm is applied to each image block and any assignments below the validation threshold are discarded. Unscorable pairs are retained: for the pairwise and assignment F1 scores, a predicted negative is used, and in retrieval the pair is given the last rank. Thus, it is possible to distinguish between failing to find a known query and correctly rejecting an unknown query and a supported non-match; since a registration failure is not resolved and may need further examination, it should not be treated as evidence of a new defect. For the assignment F1 score, which equals 2TP divided by (2TP + FP + FN), correctly accepted assignments are true positives, incorrectly accepted assignments are false positives, and known instances that are not matched are false negatives.



\subsection{Evaluation metrics and reference baselines}



The assessment of known queries makes use of Rank-1, Rank-5 and mAP, whereas pairwise decisions rely on precision, recall and the $F_1$ measure, with the oracle $F_1^{*}$ being separated from the validation-calibrated $F_1^{@v}$. The $F_1^a$ is used for assignment. In the case of identification with rejection, DIR (in the case of known accepted-correct cases) and FAR (in the case of unknown accepted cases) are used; the constrained operating point is defined by DIR@FAR$=0.1$, this point differing from the one obtained by means of a fixed threshold. Coverage is the fraction of scored pairs and is provided together with the accuracy. The references include random ranking and the best-constant pairwise $F_1$ of $2p/(1{+}p)$ (0.305 at the test prevalence of 0.180; 0.544 on the 0.48 geometric subset; 0.452 at the validation prevalence of 0.292); the full chance levels are given in Table~\ref{tab:main}. Resampling is performed under the shared capture conditions; since there are 11 test walls and 19 matchable identities we provide only point estimates, with wall-level intervals and the release of the score matrix to be considered in future work (Sec.~\ref{sec:validity}).

